% !TeX program = xelatex
% Main file. The editable text is in chapters/.
\documentclass[11pt,a4paper,oneside]{report}
\usepackage[T1]{fontenc}
\usepackage{lmodern}
\usepackage[margin=25mm,headheight=15pt]{geometry}
\usepackage{amsmath,amssymb,amsthm}
\usepackage{booktabs,enumitem,microtype}
\usepackage{xcolor,fancyhdr,needspace}
\usepackage{tikz}
\usetikzlibrary{arrows.meta}
\usepackage{hyperref}
\definecolor{ink}{HTML}{16324F}
\definecolor{muted}{HTML}{52616B}
\definecolor{accent}{HTML}{A65332}
\hypersetup{colorlinks=true,linkcolor=ink,urlcolor=ink,
 pdftitle={Algebra 1 - Chapter Notes},bookmarksnumbered=true}
\setlength{\parindent}{0pt}
\setlength{\parskip}{5pt plus 1pt}
\setlength{\emergencystretch}{2em}
\setlist{itemsep=3pt,topsep=5pt}
\newtheorem{theorem}{Theorem}[chapter]
\newtheorem{proposition}[theorem]{Proposition}
\newtheorem{lemma}[theorem]{Lemma}
\newtheorem{corollary}[theorem]{Corollary}
\theoremstyle{definition}
\newtheorem{definition}[theorem]{Definition}
\newtheorem{example}{Example}[chapter]
\theoremstyle{remark}
\newtheorem*{remark}{Remark}
\numberwithin{equation}{chapter}
\DeclareMathOperator{\spanop}{span}
\DeclareMathOperator{\rank}{rank}
\DeclareMathOperator{\im}{im}
\DeclareMathOperator{\diag}{diag}
\newcommand{\method}[1]{\par\textbf{Method.} #1\par}
\pagestyle{fancy}
\fancyhf{}
\fancyhead[L]{\small Algebra 1}
\fancyhead[R]{\small\nouppercase{\leftmark}}
\fancyfoot[C]{\thepage}
\renewcommand{\chaptermark}[1]{\markboth{Chapter \thechapter}{} }
\setcounter{tocdepth}{1}
\title{Algebra 1\\[0.4em]\large}
\author{Roger D.}
\date{2026/27}
\begin{document}
\maketitle
\pagenumbering{roman}
\chapter*{About these notes}
\addcontentsline{toc}{chapter}{About these notes}
These editable notes are based on Matthew Towers's \emph{MATH0005 Algebra 1
2026}, using the supplied revision dated 4 November 2026 (83 pages).
The earlier 6 October revision is retained as a separate source file.
This is a rewritten study companion, not an official or verbatim transcription.

The five chapters follow the source's order. Definitions, central results,
proofs and representative calculations are reorganised into a common style.
Repetitive elementary calculations are combined, while the principal methods
and all section topics are retained.

Theorem and example numbering is automatic and differs from the source.
Several important source theorem numbers are included in headings for ease
of comparison. Added methods and clarifications are identified where useful;
proofs deliberately omitted by the source are not attributed to it as if
present. The accompanying README records corrections and scope.

\textbf{Scalars.} Unless stated otherwise, $F$ means either $\mathbb R$ or
$\mathbb C$. Vectors are columns. A basis is ordered, since order determines
coordinates. The notation $N(A)$ denotes the kernel and $C(A)$ the column
space. Row-space vectors are represented as columns, following the source.

\textbf{Complex inner product convention.} We keep the source's convention
$u\cdot v=u^{\mathsf T}\overline v$, which is linear in the first argument
and conjugate-linear in the second. Other books may reverse this convention;
projection coefficients here consistently use $v\cdot u$ when projecting
$v$ onto $u$. Ordinary angle formulae are restricted to real vectors.

Edit chapter files in \texttt{chapters/}; compile this main file to combine
them with \verb|\input|. Drawings are stored separately in \texttt{figures/}.
\tableofcontents
\clearpage
\pagenumbering{arabic}
% !TeX root = ../algebra-1.tex
\chapter{Matrices and Linear Systems}\label{chap:algebra-matrices}

\section{Matrices, vectors and entrywise operations}
\begin{definition}
An $m\times n$ matrix over $F$ is a rectangular array with $m$ rows and
$n$ columns. Write $A=(a_{ij})$, where $a_{ij}$ is the entry in row $i$,
column $j$. A row vector has one row; a column vector has one column.
The set $F^n$ consists of column vectors of height $n$.
\end{definition}
For example,
\[
A=\begin{pmatrix}a&b&c\\d&e&f\end{pmatrix}
\]
is $2\times3$, its second row is $(d\ e\ f)$ and its third column is
$(c\ f)^{\mathsf T}$. The order of dimensions is always rows, then columns.
Two matrices are equal exactly when their sizes and corresponding entries agree.

Matrices of the same size add entry by entry, and scalar multiplication
multiplies every entry:
\[
(A+B)_{ij}=a_{ij}+b_{ij},\qquad(\lambda A)_{ij}=\lambda a_{ij}.
\]
For example,
\[
\begin{pmatrix}1&2\\4&5\end{pmatrix}
+\begin{pmatrix}0&1\\2&3\end{pmatrix}
=\begin{pmatrix}1&3\\6&8\end{pmatrix},\qquad
2\begin{pmatrix}1&-3\\0&1\end{pmatrix}
=\begin{pmatrix}2&-6\\0&2\end{pmatrix}.
\]
Addition is commutative and associative. For same-sized matrices,
\begin{align*}
(\alpha+\beta)A&=\alpha A+\beta A,&
\alpha(A+B)&=\alpha A+\alpha B,\\
\alpha(\beta A)&=(\alpha\beta)A.
\end{align*}
These identities follow entrywise from arithmetic in $F$.

\section{Matrix multiplication and linear combinations}
\begin{definition}[Matrix product]
If $A$ is $m\times n$ and $B$ is $n\times p$, their product is the
$m\times p$ matrix with
\[
(AB)_{ij}=\sum_{k=1}^n a_{ik}b_{kj}.
\]
The inner dimensions must agree. Each output entry is a row of $A$
multiplied into a column of $B$.
\end{definition}
\begin{example}[Sizes before arithmetic]
\[
\begin{pmatrix}a&b\\c&d\end{pmatrix}
\begin{pmatrix}e&f\\g&h\end{pmatrix}
=\begin{pmatrix}ae+bg&af+bh\\ce+dg&cf+dh\end{pmatrix}.
\]
A $2\times3$ matrix cannot multiply a $2\times1$ vector on its right.
But a $2\times3$ matrix can multiply a $3\times2$ matrix, with result
$2\times2$; for instance,
\[
\begin{pmatrix}1&2&3\\4&5&6\end{pmatrix}
\begin{pmatrix}1&2\\1&2\\1&2\end{pmatrix}
=\begin{pmatrix}6&12\\15&30\end{pmatrix}.
\]
\end{example}

\begin{definition}[Linear combination]
For same-sized matrices or vectors $v_1,\ldots,v_k$, an expression
$\sum_{j=1}^k\alpha_jv_j$, with coefficients $\alpha_j\in F$, is a linear
combination. Linear algebra studies which vectors can be produced this way
and whether their coefficients are uniquely determined.
\end{definition}

\begin{lemma}[Matrix times column; source Lemma 1.11]\label{lem:algebra-columns}
If $c_1,\ldots,c_n$ are the columns of an $m\times n$ matrix $A$ and
$x=(x_1,\ldots,x_n)^{\mathsf T}$, then
\[
Ax=x_1c_1+\cdots+x_nc_n.
\]
In particular, $Ae_j=c_j$, where $e_j$ is the standard vector with a single
$1$ in position $j$.
\end{lemma}
\begin{proof}
The $i$th entry of either side is $\sum_{j=1}^n a_{ij}x_j$.
For $x=e_j$, every coefficient except the $j$th is zero.
\end{proof}
\begin{example}
For $A=\begin{pmatrix}1&2\\3&4\end{pmatrix}$ and
$x=(1,-1)^{\mathsf T}$,
\[
Ax=\begin{pmatrix}1\\3\end{pmatrix}-\begin{pmatrix}2\\4\end{pmatrix}
=\begin{pmatrix}-1\\-1\end{pmatrix}.
\]
For the matrix with rows $(1,2,3),(4,5,6),(7,8,9)$, multiplication by
$e_2=(0,1,0)^{\mathsf T}$ selects its second column $(2,5,8)^{\mathsf T}$.
\end{example}

\begin{lemma}[Rowwise and columnwise multiplication]
If $B=[b_1\ \cdots\ b_p]$, then
$AB=[Ab_1\ \cdots\ Ab_p]$.
If $r_i$ is row $i$ of $A$, row $i$ of $AB$ is $r_iB$.
For a row $s=(s_1,\ldots,s_m)$ and rows $r_1,\ldots,r_m$ of $B$,
$sB=\sum_{i=1}^m s_ir_i$.
\end{lemma}
\begin{proof}
In each case, compare the $i,j$ entry using the product definition.
This supplies the rowwise/columnwise argument assigned as a problem in the source.
\end{proof}

\begin{example}[Multiplication encodes row and column changes]
Writing $A=\begin{pmatrix}a&b\\c&d\end{pmatrix}$,
\[
A\begin{pmatrix}0&1\\1&0\end{pmatrix}
=\begin{pmatrix}b&a\\d&c\end{pmatrix},\qquad
\begin{pmatrix}0&1\\1&1\end{pmatrix}A
=\begin{pmatrix}c&d\\a+c&b+d\end{pmatrix}.
\]
Right multiplication selects or combines columns; left multiplication
selects or combines rows.
\end{example}

\begin{proposition}[Algebra of products]
For compatible matrix sizes,
\begin{align*}
A(B+C)&=AB+AC,&(A+B)C&=AC+BC,\\
(\alpha A)B&=\alpha(AB)=A(\alpha B),&
A(BC)&=(AB)C.
\end{align*}
\end{proposition}
\begin{proof}
Distributivity and scalar compatibility follow entrywise. For associativity,
if $A$ is $m\times n$, $B$ is $n\times p$, and $C$ is $p\times q$,
\[
((AB)C)_{ij}=\sum_{\ell=1}^p\sum_{k=1}^n a_{ik}b_{k\ell}c_{\ell j}
=\sum_{k=1}^n\sum_{\ell=1}^p a_{ik}b_{k\ell}c_{\ell j}
=(A(BC))_{ij}.
\]
The finite double sums contain the same terms.
\end{proof}
Associativity allows unambiguous notation $ABC$. Commutativity generally fails:
\[
\begin{pmatrix}1&2\\3&4\end{pmatrix}
\begin{pmatrix}1&1\\1&1\end{pmatrix}
=\begin{pmatrix}3&3\\7&7\end{pmatrix},\qquad
\begin{pmatrix}1&1\\1&1\end{pmatrix}
\begin{pmatrix}1&2\\3&4\end{pmatrix}
=\begin{pmatrix}4&6\\4&6\end{pmatrix}.
\]

\section{Transpose, identity and zero matrices}
\begin{definition}[Transpose]
For $A=(a_{ij})$, its transpose $A^{\mathsf T}$ satisfies
$(A^{\mathsf T})_{ij}=a_{ji}$. It exchanges rows and columns and changes
size $m\times n$ to $n\times m$.
\end{definition}
\[
\begin{pmatrix}1&2&3\\4&5&6\end{pmatrix}^{\mathsf T}
=\begin{pmatrix}1&4\\2&5\\3&6\end{pmatrix}.
\]
The basic identities are
\begin{align*}
(A^{\mathsf T})^{\mathsf T}&=A,&
(A+B)^{\mathsf T}&=A^{\mathsf T}+B^{\mathsf T},\\
(\alpha A)^{\mathsf T}&=\alpha A^{\mathsf T},&
(AB)^{\mathsf T}&=B^{\mathsf T}A^{\mathsf T}.
\end{align*}
For the last identity, the $i,j$ entry of the left side is
$\sum_k a_{jk}b_{ki}$, equal to the $i,j$ entry of the right side.
The order reverses. Transpose alone does not conjugate complex entries.

\begin{definition}[Zero and identity]
The zero matrix $0_{m\times n}$ has every entry zero. The zero vector in
$F^n$ is $0_n$. The identity $I_n$ is $n\times n$ with $1$ on the diagonal
and $0$ elsewhere.
\end{definition}
For an $m\times n$ matrix $A$,
\[
A+0_{m\times n}=A,\qquad A I_n=A=I_mA.
\]
For square $A$, define $A^0=I_n$ and $A^k$ as the product of $k$ copies of
$A$ for $k\geq1$. Products involving zero matrices are zero with the
appropriate output size; similarly $A0_n=0_m$.

\section{Linear systems and augmented matrices}
A system of $m$ linear equations in $n$ variables has the form
\[
\sum_{j=1}^n a_{ij}x_j=b_i\quad(1\leq i\leq m).
\]
Put the coefficients into $A$, unknowns into $x$, and right sides into $b$.
The system is exactly $Ax=b$. For example,
\[
\begin{cases}2x+3y+4z=5,\\x+5z=0\end{cases}
\quad\Longleftrightarrow\quad
\begin{pmatrix}2&3&4\\1&0&5\end{pmatrix}
\begin{pmatrix}x\\y\\z\end{pmatrix}=\begin{pmatrix}5\\0\end{pmatrix}.
\]
The augmented matrix $[A\mid b]$ records coefficients and right sides together.
A solution is a vector $x$ satisfying every equation simultaneously.

\begin{definition}[Elementary row operations]
The allowed row operations are:
\begin{enumerate}
\item $R_j\leftarrow R_j+\lambda R_i$, where $i\ne j$;
\item $R_i\leftarrow\lambda R_i$, where $\lambda\ne0$;
\item $R_i\leftrightarrow R_j$.
\end{enumerate}
Their inverses subtract the same multiple, scale by $\lambda^{-1}$, and
swap the same rows, respectively.
\end{definition}
Apply operations to the \emph{entire} augmented matrix. Applying them only
to $A$ and leaving $b$ unchanged usually changes the system.

\begin{example}[Elimination with an augmented matrix]
For $3x+4y=6$, $x+2y=5$,
\begin{align*}
\left[\begin{array}{cc|c}3&4&6\\1&2&5\end{array}\right]
&\xrightarrow{R_1\leftarrow R_1-3R_2}
\left[\begin{array}{cc|c}0&-2&-9\\1&2&5\end{array}\right]\\
&\xrightarrow{R_1\leftarrow-\frac12R_1}
\left[\begin{array}{cc|c}0&1&9/2\\1&2&5\end{array}\right]\\
&\xrightarrow{R_2\leftarrow R_2-2R_1}
\left[\begin{array}{cc|c}0&1&9/2\\1&0&-4\end{array}\right]
\xrightarrow{R_1\leftrightarrow R_2}
\left[\begin{array}{cc|c}1&0&-4\\0&1&9/2\end{array}\right].
\end{align*}
Thus $x=-4$, $y=9/2$.
\end{example}

\section{Reduced row echelon form}
\begin{definition}[RREF]
The leading entry of a nonzero row is its leftmost nonzero entry.
A matrix is in reduced row echelon form (RREF) if:
\begin{enumerate}
\item every leading entry is $1$;
\item all zero rows are at the bottom;
\item a leading $1$ is the only nonzero entry in its column;
\item leading entries move strictly rightwards as one goes down nonzero rows.
\end{enumerate}
Columns containing leading entries are pivot columns.
\end{definition}
The zero matrix is in RREF. The matrix
$\begin{pmatrix}1&a&b&3\\0&0&1&-2\end{pmatrix}$ is in RREF precisely when
$b=0$. A row-echelon matrix need not have zeros \emph{above} its pivots;
RREF does.

\begin{theorem}[Existence and uniqueness of RREF]
Every matrix can be transformed by row operations into a unique RREF matrix.
\end{theorem}
\begin{proof}[Existence: elimination argument]
Scan columns from left to right. In the first available column with a nonzero
entry at or below the current pivot row, swap a nonzero entry into that row,
scale it to $1$, and eliminate all other entries in its column. Continue
below that pivot and to its right. New pivot rows have zeros in previous
pivot columns, so these operations preserve the earlier pivots.
Only finitely many rows and columns exist, and the final matrix satisfies
the RREF conditions. The source proves existence by an equivalent induction
on the number of columns. As in the source, uniqueness is stated without proof.
\end{proof}

\begin{example}[Source Example 1.34]\label{ex:algebra-rref}
Let
\[
A=\begin{pmatrix}1&2&3\\4&5&6\\7&8&9\end{pmatrix}.
\]
Subtract $4R_1$ from $R_2$ and $7R_1$ from $R_3$, then scale $R_2$ by
$-1/3$:
\[
A\sim\begin{pmatrix}1&2&3\\0&-3&-6\\0&-6&-12\end{pmatrix}
\sim\begin{pmatrix}1&2&3\\0&1&2\\0&-6&-12\end{pmatrix}
\sim\begin{pmatrix}1&0&-1\\0&1&2\\0&0&0\end{pmatrix}.
\]
The last step uses $R_1\leftarrow R_1-2R_2$ and $R_3\leftarrow R_3+6R_2$.
The pivot columns are $1$ and $2$. The symbol $\sim$ here means
row-equivalent, not equal.
\end{example}

\section{Reading the general solution}
In an augmented RREF matrix, a pivot in the last column represents an
impossible equation $0=1$. The system is then inconsistent.
Otherwise variables in coefficient pivot columns are \emph{pivot variables};
all remaining variables are \emph{free variables}.

\begin{theorem}[Free variables parameterise all solutions]
For a consistent RREF system, every assignment of scalars to the free
variables determines exactly one solution. Each nonzero row expresses its
pivot variable in terms of those free variables.
\end{theorem}
\begin{proof}
The pivot columns contain a single $1$ each. Hence each nonzero equation
has the form ``one pivot variable plus a linear combination of free variables
equals the right side''. No other equation contains that pivot variable.
Thus all pivots are uniquely determined, with no further restriction on the
free variables. Zero rows give only $0=0$.
\end{proof}

\begin{example}[Source Example 1.37]
For unknowns $(x,y,z,w)$, the augmented matrix
\[
\left[\begin{array}{cccc|c}
1&0&2&0&2\\0&1&4&0&3\\0&0&0&1&4\\0&0&0&0&0
\end{array}\right]
\]
gives $x+2z=2$, $y+4z=3$, $w=4$. Let $z=t$. Every solution is
\[
\begin{pmatrix}x\\y\\z\\w\end{pmatrix}
=\begin{pmatrix}2\\3\\0\\4\end{pmatrix}
+t\begin{pmatrix}-2\\-4\\1\\0\end{pmatrix},\qquad t\in F.
\]
Replacing the last augmented entry by $1$ produces a row $0=1$ and no solution.
\end{example}

\begin{example}[Source Example 1.39]
The system
\begin{align*}
2x_1+x_2+x_3-2x_4&=3,&x_1+x_3-2x_4&=2,\\
x_1+2x_3-4x_4&=1,&3x_1+x_2+3x_3-6x_4&=4
\end{align*}
reduces to $x_1=3$, $x_2=-2$, $x_3-2x_4=-1$. Its general solution is
\[
x=\begin{pmatrix}3\\-2\\-1\\0\end{pmatrix}
+t\begin{pmatrix}0\\0\\2\\1\end{pmatrix},\quad t\in F.
\]
Substitution into all four original equations is a useful check.
\end{example}

\begin{theorem}[More variables than equations; source Theorem 1.40]
If $A$ is $m\times n$ and $n>m$, then $Ax=0_m$ has a nonzero solution.
\end{theorem}
\begin{proof}
There are at most $m$ pivots, hence at least one free variable among the
$n$ unknowns. Give that free variable value $1$ and the other free variables
value $0$. The resulting solution is not the zero vector.
\end{proof}

\section{Inverses and elementary matrices}
\begin{definition}[Inverse]
An $n\times n$ matrix $A$ is invertible if a matrix $B$ satisfies
$AB=BA=I_n$. It is denoted $A^{-1}$ and is unique.
\end{definition}
Indeed, if $AB=I_n$ and $CA=I_n$, then
$B=(CA)B=C(AB)=C$. If $A$ is invertible, $Ax=b$ has the unique solution
$x=A^{-1}b$. The zero square matrix is not invertible, whereas $I_n$ is.
For compatible invertible square matrices,
\[
(AB)^{-1}=B^{-1}A^{-1},\qquad
(A^{\mathsf T})^{-1}=(A^{-1})^{\mathsf T}.
\]
These are checked by multiplication; inverse order reverses.

\begin{definition}[Elementary matrix]
An elementary matrix is obtained by applying one row operation to an identity
matrix. If the operation is $r$, its elementary matrix is $r(I_m)$.
\end{definition}
\begin{lemma}
For every $m\times n$ matrix $A$, performing $r$ is equivalent to left
multiplication by its elementary matrix: $r(A)=r(I_m)A$.
Elementary matrices are invertible, with inverses corresponding to the inverse
row operations. Consequently any row reduction has the form $R=EA$, where
$E$ is invertible.
\end{lemma}
\begin{proof}
Rows of $r(I_m)$ record exactly which linear combinations of the old rows
produce the new rows, so the rowwise multiplication lemma proves the identity.
Doing an inverse row operation undoes multiplication by $r(I_m)$.
For successive operations with elementary matrices $E_1,\ldots,E_k$,
$R=E_k\cdots E_1A$; the inverse product is $E_1^{-1}\cdots E_k^{-1}$.
\end{proof}
For example, $R_1\leftarrow R_1+2R_2$ for a three-row matrix is left
multiplication by
$\begin{pmatrix}1&2&0\\0&1&0\\0&0&1\end{pmatrix}$.

\begin{theorem}[Row operations preserve solutions]
If $[A'\mid b']=E[A\mid b]$ with $E$ invertible, then $Ax=b$ and
$A'x=b'$ have exactly the same solutions.
\end{theorem}
\begin{proof}
Here $A'=EA$ and $b'=Eb$. One equation follows from the other by
multiplying by $E$ or $E^{-1}$.
\end{proof}

\begin{theorem}[Little invertibility theorem; source Theorem 1.50]
For $A\in F^{n\times n}$, the following are equivalent:
$A$ is invertible; $Ax=0$ has only $x=0$; the RREF of $A$ is $I_n$.
\end{theorem}
\begin{proof}
An inverse makes $Ax=0$ imply $x=0$. If the kernel is zero, row reduction
can have no free variable, so each of the $n$ columns has a pivot and the
RREF is $I_n$. If $EA=I_n$ with $E$ invertible, then $A=E^{-1}$, so $A$
is invertible.
\end{proof}

\method{To compute an inverse, reduce $[A\mid I_n]$. If the left block can
be made $I_n$, the result is $[I_n\mid A^{-1}]$. If the left block has fewer
than $n$ pivots, $A$ is singular.}
The method works because $E[A\mid I_n]=[EA\mid E]$; when $EA=I_n$,
$E=A^{-1}$.

\begin{example}[Inverse by row reduction: source Example 1.52]
\begin{align*}
\left[\begin{array}{cc|cc}1&2&1&0\\3&4&0&1\end{array}\right]
&\sim\left[\begin{array}{cc|cc}1&2&1&0\\0&-2&-3&1\end{array}\right]\\
&\sim\left[\begin{array}{cc|cc}1&2&1&0\\0&1&3/2&-1/2\end{array}\right]
\sim\left[\begin{array}{cc|cc}1&0&-2&1\\0&1&3/2&-1/2\end{array}\right].
\end{align*}
Thus $\begin{pmatrix}1&2\\3&4\end{pmatrix}^{-1}
=\begin{pmatrix}-2&1\\3/2&-1/2\end{pmatrix}$.
For a general $2\times2$ matrix,
\[
\begin{pmatrix}a&b\\c&d\end{pmatrix}^{-1}
=\frac1{ad-bc}\begin{pmatrix}d&-b\\-c&a\end{pmatrix}
\quad\text{when }ad-bc\ne0.
\]
For larger matrices, elimination is usually more practical than an expanded
inverse formula. The determinant and cofactor formula appear in
Chapter~\ref{chap:algebra-eigenvalues}.
\end{example}

% !TeX root = ../algebra-1.tex
\chapter{Subspaces and dimension}\label{chap:algebra-subspaces}

\section{Subspaces, kernels and column spaces}

\begin{definition}
A subset $U\subseteq F^n$ is a \emph{subspace} if it is nonempty and is
closed under addition and scalar multiplication:
\[
 u,v\in U,\ a\in F \quad\Longrightarrow\quad u+v\in U,\ au\in U.
\]
We write $U\leq F^n$. More generally, $U\leq V$ means that $U$ is a
subspace contained in the subspace $V$. It is \emph{proper} if $U\ne V$.
\end{definition}

Equivalently, a nonempty subset is a subspace if $au+bv$ belongs to it
whenever $u,v$ do and $a,b\in F$.
The scalar field matters: closure must hold for every scalar in the stated
field, including negative scalars and, over $\mathbb C$, complex scalars.

\begin{proposition}
Every subspace contains the zero vector.
\end{proposition}
\begin{proof}
Choose $u\in U$, using nonemptiness. Closure under multiplication by $0$
gives $0u=0_n\in U$.
\end{proof}

\begin{example}
The sets $\{0_n\}$ and $F^n$ are subspaces. In $\mathbb R^2$, a line
through the origin, such as $\{(t,2t)^{\mathsf T}:t\in\mathbb R\}$, is a
subspace. The half-line with $t\geq0$ is not: multiplication by $-1$
takes its nonzero vectors outside the set.
A line not through the origin fails the zero-vector condition.
\end{example}

\begin{remark}
Containing zero is necessary, but does not by itself prove that a set is a
subspace. For instance, the union of the two coordinate axes contains
zero but does not contain $(1,0)^{\mathsf T}+(0,1)^{\mathsf T}$.
\end{remark}

\begin{definition}
For an $m\times n$ matrix $A$, its \emph{kernel} and \emph{column space} are
\[
 N(A)=\{v\in F^n:Av=0_m\},\qquad
 C(A)=\{Av:v\in F^n\}\subseteq F^m.
\]
The kernel describes inputs mapped to zero; the column space describes
all possible outputs.
\end{definition}

\begin{lemma}[Kernels and column spaces are subspaces]
$N(A)\leq F^n$ and $C(A)\leq F^m$.
\end{lemma}
\begin{proof}
Since $A0_n=0_m$, both sets are nonempty.
If $u,v\in N(A)$ and $a,b\in F$, then
\[
 A(au+bv)=aAu+bAv=a0_m+b0_m=0_m,
\]
so $au+bv\in N(A)$. If $y=Au$ and $z=Av$ belong to $C(A)$, then
\[
 ay+bz=aAu+bAv=A(au+bv)\in C(A).
\]
Thus each set has the required closure properties.
\end{proof}

\section{Spanning sequences and linear independence}

\begin{definition}
The \emph{span} of vectors $v_1,\ldots,v_k$ is
\[
 \spanop(v_1,\ldots,v_k)
 =\{a_1v_1+\cdots+a_kv_k:a_1,\ldots,a_k\in F\}.
\]
They \emph{span} $U$ if this set equals $U$. The span of the empty
sequence is defined to be $\{0_n\}$.
\end{definition}

\begin{proposition}
The span is a subspace and is the smallest subspace containing all the
listed vectors. The column space of $A=[\,c_1\ \cdots\ c_n\,]$ is
$\spanop(c_1,\ldots,c_n)$.
\end{proposition}
\begin{proof}
Zero is a linear combination with all coefficients zero. Sums and scalar
multiples of linear combinations are again linear combinations.
Any subspace containing the listed vectors contains all their linear
combinations, which proves minimality. Finally,
$Ax=x_1c_1+\cdots+x_nc_n$, so the possible outputs are exactly their span.
\end{proof}

\begin{example}
The $xy$-plane in $\mathbb R^3$ is spanned by $e_1,e_2$, and also by
$(1,1,0)^{\mathsf T},(1,-1,0)^{\mathsf T}$, since
\[
 (a,b,0)^{\mathsf T}
 =\frac{a+b}{2}(1,1,0)^{\mathsf T}
  +\frac{a-b}{2}(1,-1,0)^{\mathsf T}.
\]
The sequence $e_1,e_2,e_3$ spans all of $\mathbb R^3$, so it does not
span this plane: its span is too large.
\end{example}

\begin{example}[Testing a span by solving a system]
Let
\[
 u=\begin{pmatrix}1\\2\\3\end{pmatrix},\quad
 v=\begin{pmatrix}4\\5\\6\end{pmatrix},\quad
 w=\begin{pmatrix}7\\8\\9\end{pmatrix}.
\]
To decide whether $(a,b,c)^{\mathsf T}$ belongs to their span, solve the
system with coefficient matrix $[\,u\ v\ w\,]$. Row reduction gives
\[
 \left[\begin{array}{ccc|c}
 1&4&7&a\\2&5&8&b\\3&6&9&c
 \end{array}\right]
 \ \longrightarrow\
 \left[\begin{array}{ccc|c}
 1&0&-1&(4b-5a)/3\\
 0&1&2&(2a-b)/3\\
 0&0&0&a-2b+c
 \end{array}\right].
\]
Consequently,
\[
 \spanop(u,v,w)
 =\{(a,b,c)^{\mathsf T}:a-2b+c=0\}.
\]
This is a plane, not all of $\mathbb R^3$.
\end{example}

\begin{definition}
A sequence $v_1,\ldots,v_k$ is \emph{linearly independent} if
\[
 a_1v_1+\cdots+a_kv_k=0_n
 \quad\Longrightarrow\quad a_1=\cdots=a_k=0.
\]
Otherwise it is \emph{linearly dependent}. The empty sequence is independent.
\end{definition}

A sequence containing zero is dependent. A sequence with a repeated vector
is dependent. More generally, dependence means that at least one listed
vector can be expressed as a linear combination of the others.

\begin{example}
The vectors $(1,0)^{\mathsf T},(0,1)^{\mathsf T},(1,1)^{\mathsf T}$
are dependent because the first two minus the third give zero.
By contrast, $(1,1)^{\mathsf T},(1,-1)^{\mathsf T}$ are independent:
\[
 a(1,1)^{\mathsf T}+b(1,-1)^{\mathsf T}=0
 \quad\Longrightarrow\quad a+b=a-b=0
 \quad\Longrightarrow\quad a=b=0.
\]
For the preceding three-dimensional example, $u-2v+w=0$, so $u,v,w$
are dependent.
\end{example}

\method{Put the vectors in the columns of a matrix $A$. They are
independent exactly when $Ax=0$ has only the zero solution.
To test whether they span $F^m$, ask whether $Ax=b$ is solvable for
every $b\in F^m$. These are different questions.}

\section{Bases and dimension}

\begin{definition}
A \emph{basis} of $U$ is an ordered sequence of vectors that is both
linearly independent and spans $U$.
\end{definition}

Spanning ensures that every vector has a representation in the basis.
Independence ensures that this representation is unique: subtract two
representations and apply the definition of independence.
The standard basis of $F^n$ is $e_1,\ldots,e_n$.

\begin{lemma}[Steinitz length bound; source Lemma 2.15]
Suppose $s_1,\ldots,s_m$ span $U$. Any sequence of more than $m$ vectors
in $U$ is linearly dependent.
\end{lemma}
\begin{proof}
Let $v_1,\ldots,v_k\in U$ with $k>m$. Write
$v_j=\sum_{i=1}^m a_{ij}s_i$, and form the $m\times k$ matrix
$A=(a_{ij})$. The homogeneous system $Ax=0$ has more unknowns than
equations, so it has a nonzero solution $x$. Then
\[
 \sum_{j=1}^k x_jv_j
 =\sum_{i=1}^m\left(\sum_{j=1}^k a_{ij}x_j\right)s_i=0.
\]
Since some $x_j$ is nonzero, this is a nontrivial dependence.
The case $m=0$ simply says that every vector of $U=\{0\}$ is zero.
\end{proof}

\begin{corollary}
Any $n+1$ vectors in $F^n$ are dependent.
\end{corollary}
\begin{proof}
Apply the lemma to the $n$ standard basis vectors.
\end{proof}

\begin{theorem}[Size of a basis; source Theorem 2.17]
Any two bases of a subspace have the same number of vectors.
\end{theorem}
\begin{proof}
Let one basis have $r$ vectors and the other $s$.
The first spans the subspace and the second is independent, so the
Steinitz bound gives $s\leq r$. Reversing the roles gives $r\leq s$.
Thus $r=s$.
\end{proof}

\begin{lemma}[Extending an independent sequence]
If $u_1,\ldots,u_k$ are independent and
$v\notin\spanop(u_1,\ldots,u_k)$, then $u_1,\ldots,u_k,v$ are independent.
\end{lemma}
\begin{proof}
If $\sum_i a_i u_i+bv=0$ and $b\ne0$, then
$v=-b^{-1}\sum_i a_i u_i$, contrary to the assumption.
Thus $b=0$, after which independence of the $u_i$ gives all $a_i=0$.
\end{proof}

\begin{theorem}[Basis extension]
Every independent sequence in a subspace $U\leq F^n$ can be extended to
a basis of $U$. In particular, every subspace has a basis.
\end{theorem}
\begin{proof}
If the current sequence does not span $U$, choose a vector of $U$ outside
its span and append it. The lemma preserves independence.
This process cannot produce more than $n$ independent vectors, so it
must stop at a spanning sequence. Starting with the empty sequence gives
existence, including the empty basis of $\{0_n\}$.
\end{proof}

\begin{definition}
The \emph{dimension} of $U$, written $\dim U$, is the number of vectors
in any basis of $U$. The size theorem makes this well-defined.
\end{definition}

\begin{example}
$\dim F^n=n$ and $\dim\{0_n\}=0$.
The plane $a+b+c=0$ has basis
$(1,0,-1)^{\mathsf T},(0,1,-1)^{\mathsf T}$, because
$(a,b,c)^{\mathsf T}=a(1,0,-1)^{\mathsf T}
+b(0,1,-1)^{\mathsf T}$ whenever $c=-a-b$.
Its dimension is two.
\end{example}

\begin{proposition}
If $U\leq V\leq F^n$, then $\dim U\leq\dim V$.
Equality holds if and only if $U=V$.
\end{proposition}
\begin{proof}
Extend a basis of $U$ to a basis of $V$. This can only increase its
length. If the dimensions are equal, no extra vectors are needed, so
the original basis already spans $V$ and $U=V$.
\end{proof}

\begin{corollary}
In a $d$-dimensional subspace, any $d$ independent vectors form a basis.
Also, any $d$ vectors that span it form a basis.
\end{corollary}
\begin{proof}
For the first assertion, basis extension leaves no room for extra vectors.
For the second, if a spanning sequence is dependent, delete a vector
expressible using the others. Repeating gives a basis of length less
than $d$, a contradiction.
\end{proof}

\section{Sums and intersections}

\begin{definition}
For subspaces $S,T\leq F^n$, define
\[
 S\cap T=\{v:v\in S\text{ and }v\in T\},\qquad
 S+T=\{s+t:s\in S,\ t\in T\}.
\]
\end{definition}

Both are subspaces. For the intersection, a linear combination stays in
each of the two sets. For the sum,
$a(s_1+t_1)+b(s_2+t_2)=(as_1+bs_2)+(at_1+bt_2)\in S+T$.
Both sets contain zero.
The sum is the span of the combined spanning sequences of $S$ and $T$.
It usually differs from the union $S\cup T$.

\begin{example}
Let $S=\spanop((2,-1,0)^{\mathsf T})$ and
$T=\spanop((2,-1,2)^{\mathsf T})$ in $\mathbb R^3$.
The two generating vectors are independent, so their sum is a plane.
If $a(2,-1,0)^{\mathsf T}=b(2,-1,2)^{\mathsf T}$, the third coordinate
gives $b=0$, and then $a=0$. Thus $S\cap T=\{0_3\}$.
\end{example}

\begin{theorem}[Dimension formula; source Theorem 2.28]
\[
 \boxed{\dim(S+T)=\dim S+\dim T-\dim(S\cap T).}
\]
\end{theorem}
\begin{proof}
Choose a basis $v_1,\ldots,v_k$ of $S\cap T$.
Extend it to a basis $v_1,\ldots,v_k,s_1,\ldots,s_r$ of $S$, and
to a basis $v_1,\ldots,v_k,t_1,\ldots,t_q$ of $T$.
The combined sequence $v_1,\ldots,v_k,s_1,\ldots,s_r,t_1,\ldots,t_q$
spans $S+T$.

For independence, suppose
$\sum_i a_iv_i+\sum_j b_js_j+\sum_\ell c_\ell t_\ell=0$.
Then $\sum_\ell c_\ell t_\ell\in S\cap T$, so it is some
$\sum_i d_iv_i$. Independence of the chosen basis of $T$ forces
all $c_\ell$ and $d_i$ to be zero.
Independence of the basis of $S$ then gives all $a_i=b_j=0$.
Thus the combined sequence is a basis of length $k+r+q$, which equals
$(k+r)+(k+q)-k$.
\end{proof}

For example, two distinct planes through the origin in $\mathbb R^3$
intersect in a line and have sum $\mathbb R^3$.
If $S\cap T=\{0\}$, the sum is called a \emph{direct sum}, written
$S\oplus T$; the representation $s+t$ is then unique.

\section{Computing bases from a matrix}

\subsection{A basis for the kernel}

For an RREF matrix, assign each free variable a parameter.
Set one free variable to $1$ and the others to $0$, then solve for the
pivot variables. The resulting vectors are the \emph{fundamental solutions}.

\begin{example}[Source Example 2.30]
Suppose $A$ has reduced row echelon form
\[
 R=\begin{pmatrix}0&1&2&0&3\\0&0&0&1&4\\0&0&0&0&0\end{pmatrix}.
\]
The equations are $x_2=-2x_3-3x_5$ and $x_4=-4x_5$.
Variables $x_1,x_3,x_5$ are free, so
\[
 N(A)=N(R)=\spanop\left(
 \begin{pmatrix}1\\0\\0\\0\\0\end{pmatrix},
 \begin{pmatrix}0\\-2\\1\\0\\0\end{pmatrix},
 \begin{pmatrix}0\\-3\\0\\-4\\1\end{pmatrix}\right).
\]
These vectors are independent: their first, third and fifth coordinates
recover their coefficients. They therefore form a basis, and $\dim N(A)=3$.
\end{example}

\begin{proposition}
The fundamental solutions form a basis of $N(A)$. Hence $\dim N(A)$
is the number of free variables.
\end{proposition}
\begin{proof}
Every solution is the linear combination whose coefficients are its
free-variable values, so the fundamental solutions span.
At the free-variable positions they have the columns of an identity
matrix, so a zero linear combination must have every coefficient zero.
\end{proof}

\subsection{A basis for the column space}

\begin{definition}
The \emph{column rank}, or simply \emph{rank}, is
$\rank A=\dim C(A)$.
\end{definition}

\begin{theorem}
If the pivot columns of the RREF of $A$ have indices
$j_1,\ldots,j_r$, then columns $j_1,\ldots,j_r$ of the
\textbf{original matrix $A$} form a basis of $C(A)$.
\end{theorem}
\begin{proof}
Write $R=EA$, where $E$ is the invertible product of the elementary
matrices used in row reduction. For every coefficient vector $x$,
$Ax=0$ if and only if $Rx=0$. Thus the columns of $A$ and $R$ have
exactly the same linear relations.

In $R$, the pivot columns are the first $r$ standard coordinate vectors.
Every other column is their linear combination.
The identical relations hold in $A$, so its corresponding columns span
$C(A)$. Any dependence among these columns would give a dependence among
the pivot columns of $R$, which are independent.
\end{proof}

\begin{remark}
Row operations preserve the kernel and the relations between columns,
but can change the column space itself. For example,
\[
 A=\begin{pmatrix}1&0\\1&0\end{pmatrix}
 \longrightarrow R=\begin{pmatrix}1&0\\0&0\end{pmatrix}.
\]
Here $C(A)=\spanop((1,1)^{\mathsf T})$ and
$C(R)=\spanop((1,0)^{\mathsf T})$ are different.
\end{remark}

\begin{example}
The matrix in Example~\ref{ex:algebra-rref} has pivots in columns one and
two. Therefore a column-space basis is
$(1,4,7)^{\mathsf T},(2,5,8)^{\mathsf T}$, and its rank is two.
\end{example}

\begin{theorem}[Rank--nullity; source Theorem 2.39]
For every $m\times n$ matrix $A$,
\[
 \boxed{\dim N(A)+\rank A=n.}
\]
\end{theorem}
\begin{proof}
There are $n$ columns. Each is either a pivot column or a free-variable
column. The two numbers are $\rank A$ and $\dim N(A)$ respectively.
\end{proof}

\subsection{A basis for the row space}

Representing rows as column vectors, the \emph{row space} of $A$ is
$C(A^{\mathsf T})\leq F^n$.

\begin{lemma}
If $F_0$ is an invertible $n\times n$ matrix, then $C(AF_0)=C(A)$.
\end{lemma}
\begin{proof}
Every $AF_0x$ is in $C(A)$. Conversely, every $Ay$ equals
$AF_0(F_0^{-1}y)$, so it lies in $C(AF_0)$.
\end{proof}

Consequently, row operations preserve row space:
if $R=EA$, then $R^{\mathsf T}=A^{\mathsf T}E^{\mathsf T}$ and the
lemma applies. The nonzero rows of $R$, transposed, form a row-space
basis. They span by definition, and their pivot positions prove
independence.

\begin{theorem}[Equality of row and column rank]
\[
 \dim C(A^{\mathsf T})=\dim C(A)=\rank A.
\]
\end{theorem}
\begin{proof}
Both dimensions equal the number of pivots in an echelon form.
\end{proof}

Equal dimensions do not mean equal spaces. For
$A=\begin{pmatrix}1&1\\2&2\end{pmatrix}$, the column space is
$\spanop((1,2)^{\mathsf T})$, while the row space is
$\spanop((1,1)^{\mathsf T})$.

\begin{center}
\begin{tabular}{lll}
\toprule
Space & Ambient space & Basis obtained from RREF $R$\\
\midrule
$N(A)$ & $F^n$ & Fundamental solutions of $Rx=0$\\
$C(A)$ & $F^m$ & Original columns at pivot indices\\
$C(A^{\mathsf T})$ & $F^n$ & Nonzero rows of $R$, transposed\\
\bottomrule
\end{tabular}
\end{center}

\section{Equivalent tests for invertibility}

\begin{theorem}[Big invertibility theorem; source Theorem 2.40]
For a square $n\times n$ matrix $A$, the following are equivalent:
\begin{enumerate}
\item $A$ is invertible.
\item $A^{\mathsf T}$ is invertible.
\item The RREF of $A$ is $I_n$.
\item $N(A)=\{0_n\}$.
\item The columns of $A$ are independent.
\item $\rank A=n$.
\item The columns of $A$ form a basis of $F^n$.
\item $C(A)=F^n$.
\item For every $b\in F^n$, the equation $Ax=b$ has a unique solution.
\end{enumerate}
\end{theorem}
\begin{proof}
The inverse criterion from Chapter~\ref{chap:algebra-matrices} equates
invertibility with RREF $I_n$.
For a square matrix, this means there are $n$ pivots, equivalently no
free variables, equivalently a trivial kernel.
A trivial kernel is precisely independence of the columns.
The column-space basis theorem gives rank equal to the number of pivots.
In dimension $n$, $n$ independent columns form a basis, equivalently they
span $F^n$. In that basis, every $b$ has a unique coefficient vector
$x$, which is the unique solution of $Ax=b$.
Conversely, unique solvability at $b=0$ gives a trivial kernel.
Finally, if $A$ is invertible, then
$(A^{\mathsf T})^{-1}=(A^{-1})^{\mathsf T}$; applying the same argument
to $A^{\mathsf T}$ gives the reverse implication.
\end{proof}

The determinant criterion $\det A\ne0$ will be added in
Chapter~\ref{chap:algebra-eigenvalues}.
For a rectangular matrix, keep the distinctions:
independent columns mean rank $n$, whereas spanning the codomain $F^m$
means rank $m$. They coincide only when $m=n$.

% !TeX root = ../algebra-1.tex
\chapter{Linear maps and coordinates}\label{chap:algebra-linear-maps}

\section{Matrices as linear maps}

An $m\times n$ matrix defines a function
\[
 T_A:F^n\longrightarrow F^m,\qquad T_A(v)=Av.
\]
Its input has $n$ coordinates and its output has $m$ coordinates.
For example, $\begin{pmatrix}-1&0\\0&1\end{pmatrix}$ reflects
$\mathbb R^2$ in the $y$-axis.

\begin{definition}
A map $T:F^n\to F^m$ is \emph{linear} if
\[
 T(u+v)=T(u)+T(v),\qquad T(av)=aT(v)
\]
for all $u,v\in F^n$ and $a\in F$.
Equivalently, $T(au+bv)=aT(u)+bT(v)$ for all $a,b,u,v$.
\end{definition}

Matrix multiplication shows that $T_A$ is linear. Repeated application of
linearity gives
$T(\sum_i a_iv_i)=\sum_i a_iT(v_i)$.
Every linear map sends zero to zero, since $T(0v)=0T(v)=0_m$.
A nonzero translation $v\mapsto v+c$ is therefore not linear.

\begin{proposition}[Composition]
If $B$ is $n\times p$ and $A$ is $m\times n$, then
\[
 T_A\circ T_B=T_{AB}.
\]
\end{proposition}
\begin{proof}
For $v\in F^p$, $T_A(T_B(v))=A(Bv)=(AB)v$.
\end{proof}

The rightmost map acts first. An invertible matrix represents a linear
bijection, and its inverse map is represented by the inverse matrix.

\begin{theorem}[Every linear map has a matrix; source Theorem 3.4]
For each linear $T:F^n\to F^m$, there is a unique $m\times n$ matrix $A$
such that $T(v)=Av$ for every $v$. Its columns are
\[
 \boxed{A=[\,T(e_1)\ \cdots\ T(e_n)\,].}
\]
\end{theorem}
\begin{proof}
Write $v=\sum_j v_je_j$. Linearity gives
$T(v)=\sum_j v_jT(e_j)$, exactly the column formula for $Av$.
Conversely, a representing matrix must have column $j$ equal to
$Ae_j=T(e_j)$, so it is unique.
\end{proof}

\begin{example}
The map $T(a,b,c)=a+b+c$ has matrix $(1\ 1\ 1)$.
The map $S(a,b)=(a,b,a,b)^{\mathsf T}$ has matrix
\[
 \begin{pmatrix}1&0\\0&1\\1&0\\0&1\end{pmatrix}.
\]
Finding images of the standard basis is often quicker than guessing
matrix entries.
\end{example}

\section{Kernel, image and invertibility of a map}

\begin{definition}
The \emph{kernel} and \emph{image} of a linear map are
\[
 \ker T=\{v:T(v)=0_m\},\qquad
 \im T=\{T(v):v\in F^n\}.
\]
If $A$ represents $T$, these are $N(A)$ and $C(A)$.
Its rank is $\dim(\im T)$.
\end{definition}

They are subspaces by the same closure proofs as for matrices.
Rank--nullity becomes
\[
 \dim(\ker T)+\dim(\im T)=n.
\]

\begin{example}
For $T(a,b,c)=a+b+c$, the kernel is the plane $a+b+c=0$,
with basis $(1,0,-1)^{\mathsf T},(0,1,-1)^{\mathsf T}$.
The image is all of $F$, since $T(t,0,0)=t$.
Thus the kernel dimension is two and the rank is one.

For $S(a,b)=(a,b,a,b)^{\mathsf T}$, the kernel is $\{0_2\}$ and
\[
 \im S=\spanop\left(
 (1,0,1,0)^{\mathsf T},(0,1,0,1)^{\mathsf T}\right).
\]
This map is injective, but is not surjective onto $F^4$.
\end{example}

\begin{proposition}
A linear map is injective if and only if its kernel is $\{0_n\}$.
It is surjective if and only if its image is the whole codomain.
\end{proposition}
\begin{proof}
If it is injective, $T(v)=0_m=T(0_n)$ implies $v=0_n$.
Conversely, if its kernel is trivial and $T(u)=T(v)$, then
$T(u-v)=0_m$ gives $u=v$.
The second assertion is the definition of surjectivity.
\end{proof}

Thus for $T:F^n\to F^m$, injectivity means rank $n$ and surjectivity
means rank $m$. Rank--nullity implies that an injective map requires
$n\leq m$ and a surjective map requires $n\geq m$.
When $m=n$, injectivity and surjectivity are equivalent, and both mean
that the representing matrix is invertible.

\section{Coordinates in an ordered basis}

\begin{definition}
For an ordered basis $\mathcal B=(b_1,\ldots,b_n)$ of $F^n$, the
\emph{coordinate vector} of $v$ is
\[
 [v]_{\mathcal B}=\begin{pmatrix}c_1\\\vdots\\c_n\end{pmatrix}
 \quad\text{where}\quad v=c_1b_1+\cdots+c_nb_n.
\]
\end{definition}

Let $P=[\,b_1\ \cdots\ b_n\,]$. Its columns form a basis, so $P$ is
invertible and
\[
 \boxed{v=P[v]_{\mathcal B},\qquad [v]_{\mathcal B}=P^{-1}v.}
\]
The vector $v$ is unchanged; these equations describe the same vector
using different coordinate systems.
The direction is worth remembering:
\[
 \mathcal B\text{-coordinates}
 \ \xrightarrow{\ P\ }\ \text{standard coordinates},
 \qquad
 \text{standard coordinates}
 \ \xrightarrow{\ P^{-1}\ }\ \mathcal B\text{-coordinates}.
\]

\begin{example}
Take $b_1=(1,1)^{\mathsf T}$ and $b_2=(1,-1)^{\mathsf T}$.
Then
\[
 P=\begin{pmatrix}1&1\\1&-1\end{pmatrix},\qquad
 P^{-1}=\frac12\begin{pmatrix}1&1\\1&-1\end{pmatrix}.
\]
For $v=(0,1)^{\mathsf T}$, $[v]_{\mathcal B}=(1/2,-1/2)^{\mathsf T}$:
\[
 \begin{pmatrix}0\\1\end{pmatrix}
 =\frac12\begin{pmatrix}1\\1\end{pmatrix}
  -\frac12\begin{pmatrix}1\\-1\end{pmatrix}.
\]
\end{example}

\section{Changing the matrix of a map}

\begin{theorem}[Change of basis; source Theorem 3.14]
Suppose $T:F^n\to F^m$ has standard matrix $A$.
Let $\mathcal B$ be an input basis with column matrix $P$, and
$\mathcal C$ an output basis with column matrix $Q$.
The matrix taking $\mathcal B$-coordinates of the input to
$\mathcal C$-coordinates of the output is
\[
 \boxed{[T]_{\mathcal C\leftarrow\mathcal B}=Q^{-1}AP.}
\]
\end{theorem}
\begin{proof}
If $x=[v]_{\mathcal B}$, then $v=Px$.
The standard output is $APx$, and its $\mathcal C$-coordinates are
$Q^{-1}APx$. Equivalently, column $j$ of the new matrix is
$[T(b_j)]_{\mathcal C}$.
\end{proof}

When $T:F^n\to F^n$ uses the same basis for its input and output, this
becomes $P^{-1}AP$.
Matrices related by $B=P^{-1}AP$ are called \emph{similar}.
Similarity describes one operator in different bases.
Row reduction instead multiplies on the left by elementary matrices;
it serves a different purpose and is not generally similarity.

\begin{example}[A basis adapted to the map]
Let $A=\begin{pmatrix}2&1\\1&2\end{pmatrix}$ and use the preceding basis.
Since
\[
 A\begin{pmatrix}1\\1\end{pmatrix}
 =3\begin{pmatrix}1\\1\end{pmatrix},\qquad
 A\begin{pmatrix}1\\-1\end{pmatrix}
 =\begin{pmatrix}1\\-1\end{pmatrix},
\]
we obtain
\[
 P^{-1}AP=\begin{pmatrix}3&0\\0&1\end{pmatrix}.
\]
In these coordinates the map stretches one coordinate by $3$ and leaves
the other unchanged.
\end{example}

\section{Diagonalisation and matrix powers}

If $D=\diag(d_1,\ldots,d_n)$, then
$D^k=\diag(d_1^k,\ldots,d_n^k)$ for positive integers $k$.
When $P^{-1}AP=D$, equivalently $A=PDP^{-1}$, cancellation of the middle
factors gives
\[
 \boxed{A^k=PD^kP^{-1}\quad(k\geq0).}
\]
Here $D^0=I_n$. Negative powers are also valid when all $d_i$ are
nonzero, so that $A$ is invertible.

\begin{example}
For $A=\begin{pmatrix}3&-1\\-1&3\end{pmatrix}$ and the same $P$,
\[
 P^{-1}AP=\diag(2,4).
\]
Therefore, for $k\geq0$,
\[
 A^k=\frac12
 \begin{pmatrix}
 2^k+4^k&2^k-4^k\\
 2^k-4^k&2^k+4^k
 \end{pmatrix}.
\]
The directions $(1,1)^{\mathsf T}$ and $(1,-1)^{\mathsf T}$ are
scaled by $2^k$ and $4^k$ respectively.
\end{example}

\subsection{Application: the Fibonacci sequence}

Define $f_0=0$, $f_1=1$ and $f_{k+2}=f_{k+1}+f_k$.
With
\[
 M=\begin{pmatrix}1&1\\1&0\end{pmatrix},
 \qquad
 \begin{pmatrix}f_{k+1}\\f_k\end{pmatrix}
 =M^k\begin{pmatrix}1\\0\end{pmatrix},
\]
the recurrence becomes a matrix power.
To find invariant directions, calculate
$M(\lambda,1)^{\mathsf T}=(\lambda+1,\lambda)^{\mathsf T}$.
This equals $\lambda(\lambda,1)^{\mathsf T}$ when
$\lambda^2=\lambda+1$.
The two roots are
\[
 \phi_+=\frac{1+\sqrt5}{2},\qquad
 \phi_-=\frac{1-\sqrt5}{2}.
\]
Set
\[
 P=\begin{pmatrix}\phi_+&\phi_-\\1&1\end{pmatrix},\qquad
 P^{-1}=\frac1{\sqrt5}
 \begin{pmatrix}1&-\phi_-\\-1&\phi_+\end{pmatrix}.
\]
The invariant-direction calculations give
$MP=P\diag(\phi_+,\phi_-)$.
Hence
\begin{align*}
 \begin{pmatrix}f_{k+1}\\f_k\end{pmatrix}
 &=P\begin{pmatrix}\phi_+^k&0\\0&\phi_-^k\end{pmatrix}
       \frac1{\sqrt5}\begin{pmatrix}1\\-1\end{pmatrix}\\
 &=\frac1{\sqrt5}
   \begin{pmatrix}\phi_+^{k+1}-\phi_-^{k+1}\\
                   \phi_+^k-\phi_-^k\end{pmatrix}.
\end{align*}
The second coordinate is \emph{Binet's formula}:
\[
 \boxed{f_k=\frac{\phi_+^k-\phi_-^k}{\sqrt5}.}
\]
The invariant directions used here are eigenvectors, the main topic of
the next chapter.

% !TeX root = ../algebra-1.tex
\chapter{Eigenvalues, diagonalisation and determinants}
\label{chap:algebra-eigenvalues}

\section{Eigenvectors and eigenspaces}

\begin{definition}
Let $A$ be an $n\times n$ matrix over $F$.
A \emph{nonzero} vector $v\in F^n$ is an \emph{eigenvector} with
\emph{eigenvalue} $\lambda\in F$ if
\[
 Av=\lambda v.
\]
The \emph{eigenspace} for $\lambda$ is
\[
 E_\lambda(A)=N(A-\lambda I_n).
\]
\end{definition}

An eigenvector is a direction preserved by the map: its image is a
scalar multiple of itself. Zero is never an eigenvector, although zero
can be an eigenvalue.
The eigenspace is a subspace and includes zero; its nonzero vectors are
exactly the eigenvectors for that eigenvalue.

\begin{example}[The all-ones matrix]
Let $J$ be the $3\times3$ matrix all of whose entries are $1$.
For $v=(a,b,c)^{\mathsf T}$,
$Jv=(a+b+c)(1,1,1)^{\mathsf T}$.
It follows that
\begin{align*}
 E_0(J)&=\{(a,b,c)^{\mathsf T}:a+b+c=0\}\\
 &=\spanop\bigl((-1,1,0)^{\mathsf T},(-1,0,1)^{\mathsf T}\bigr),\\
 E_3(J)&=\spanop\bigl((1,1,1)^{\mathsf T}\bigr).
\end{align*}
These are the only eigenvalues: if $Jv=\lambda v$ and $\lambda\ne0$,
then $v$ is a multiple of $(1,1,1)^{\mathsf T}$, which forces $\lambda=3$.
\end{example}

\begin{example}[The field matters]
For $A=\begin{pmatrix}0&1\\-1&0\end{pmatrix}$, we have $A^2=-I_2$.
If $Av=\lambda v$ with $v\ne0$, then $\lambda^2=-1$, so there are
no real eigenvalues or real eigenvectors.
Over $\mathbb C$ there are two eigenvalues:
\[
 A\begin{pmatrix}-i\\1\end{pmatrix}
 =i\begin{pmatrix}-i\\1\end{pmatrix},\qquad
 A\begin{pmatrix}i\\1\end{pmatrix}
 =-i\begin{pmatrix}i\\1\end{pmatrix}.
\]
\end{example}

\section{When is a matrix diagonalisable?}

\begin{definition}
$A$ is \emph{diagonalisable over $F$} if there is an invertible matrix
$P$ over $F$ such that $P^{-1}AP$ is diagonal.
\end{definition}

\begin{theorem}[Eigenvector basis criterion]
$A$ is diagonalisable over $F$ if and only if $F^n$ has a basis of
eigenvectors of $A$.
If $P=[\,v_1\ \cdots\ v_n\,]$ consists of such a basis and
$Av_j=\lambda_jv_j$, then
\[
 P^{-1}AP=\diag(\lambda_1,\ldots,\lambda_n).
\]
\end{theorem}
\begin{proof}
The equality $AP=PD$ is an equality of columns:
its $j$th column says $Av_j=\lambda_jv_j$.
Invertibility of $P$ means that its columns form a basis.
Multiplying by $P^{-1}$ proves both directions.
\end{proof}

The order of the eigenvalues on the diagonal must match the order of the
eigenvectors in $P$.
The preceding rotation matrix is diagonalisable over $\mathbb C$, but not
over $\mathbb R$.

\begin{example}[A matrix that is not diagonalisable]
For
\[
 A=\begin{pmatrix}1&1\\0&1\end{pmatrix},
 \qquad E_1(A)=\spanop((1,0)^{\mathsf T}).
\]
Indeed, $Av=\lambda v$ first gives
$(1-\lambda)y=0$ and $(1-\lambda)x+y=0$.
A nonzero solution forces $\lambda=1$, and then $y=0$.
Thus there are not enough independent eigenvectors to form a basis.
\end{example}

\begin{lemma}[Independence of distinct eigenspaces]
Let $\lambda_1,\ldots,\lambda_r$ be distinct eigenvalues and
$v_j\in E_{\lambda_j}(A)$.
If $v_1+\cdots+v_r=0$, then every $v_j=0$.
\end{lemma}
\begin{proof}
Use induction on $r$. The case $r=1$ is immediate.
Apply $A-\lambda_r I$ to the given relation:
\[
 \sum_{j=1}^{r-1}(\lambda_j-\lambda_r)v_j=0.
\]
The induction hypothesis implies
$(\lambda_j-\lambda_r)v_j=0$ for $j<r$.
Distinctness makes these scalar factors nonzero, so $v_j=0$ for $j<r$;
the original relation then gives $v_r=0$.
\end{proof}

\begin{corollary}
Eigenvectors with distinct eigenvalues are linearly independent.
An $n\times n$ matrix has at most $n$ distinct eigenvalues.
If it has $n$ distinct eigenvalues in $F$, it is diagonalisable over $F$.
\end{corollary}
\begin{proof}
Apply the lemma to $a_jv_j$ in a proposed linear relation.
Nonzero $v_j$ then force all coefficients to vanish.
There cannot be more than $n$ independent vectors in $F^n$, while
$n$ independent vectors form a basis.
\end{proof}

\begin{theorem}[Dimension criterion]
If $\lambda_1,\ldots,\lambda_r$ are all the eigenvalues in $F$, then
\[
 A\text{ is diagonalisable over }F
 \quad\Longleftrightarrow\quad
 \sum_{j=1}^r\dim E_{\lambda_j}(A)=n.
\]
\end{theorem}
\begin{proof}
Choose a basis in each eigenspace. The preceding lemma shows that their
combined sequence is independent: group a linear relation according to
eigenvalue, then use independence within each eigenspace.
The combined sequence is an eigenvector basis exactly when it has
length $n$. If an eigenvector basis already exists, each of its vectors
belongs to one of these eigenspaces, so their sum is all of $F^n$.
\end{proof}

Distinct eigenvalues are sufficient, but not necessary.
For example, $I_n$ is already diagonal although it has only one
eigenvalue.

\section{Determinants}

The source uses the determinant and its properties without constructing
it from first principles. We follow that scope here.
For a square matrix, $\det A$ is a scalar satisfying, in particular,
\[
 \det I_n=1,\qquad \det(AB)=\det A\,\det B,\qquad
 A\text{ invertible}\ \Longleftrightarrow\ \det A\ne0.
\]
The last condition extends the invertibility theorem in
Chapter~\ref{chap:algebra-subspaces}.
The determinant is not generally additive: $\det(A+B)$ need not equal
$\det A+\det B$.

For small matrices,
\[
 \det\begin{pmatrix}a&b\\c&d\end{pmatrix}=ad-bc,
\]
and expansion along the first row gives
\[
 \det\begin{pmatrix}a&b&c\\d&e&f\\g&h&i\end{pmatrix}
 =a(ei-fh)-b(di-fg)+c(dh-eg).
\]
The signs in this expansion are $+,-,+$.

\subsection{Computing determinants by elimination}

The elementary row operations affect determinants as follows.
\begin{center}
\begin{tabular}{ll}
\toprule
Operation & Effect on determinant\\
\midrule
Swap two rows & Multiply by $-1$\\
Multiply one row by $c$ & Multiply by $c$\\
Add a multiple of one row to another & No change\\
\bottomrule
\end{tabular}
\end{center}
For an upper or lower triangular matrix,
\[
 \boxed{\det A=a_{11}a_{22}\cdots a_{nn}.}
\]
For an upper triangular matrix with nonzero diagonal, eliminate entries
above the diagonal using row additions to obtain a diagonal matrix;
its determinant is the product of its diagonal entries.
If a diagonal entry is zero, such a triangular matrix is singular
(its homogeneous system has a nonzero solution by backward solving),
so its determinant and its diagonal product are both zero.
The lower triangular case follows similarly.

\method{Reduce to triangular form, keeping track of every row swap and
row scaling. Multiply the diagonal entries, then undo the recorded
factors. Using only row additions avoids scaling factors.}

\begin{example}
Using row additions alone,
\[
 \begin{pmatrix}1&2&3\\4&5&6\\7&8&10\end{pmatrix}
 \longrightarrow
 \begin{pmatrix}1&2&3\\0&-3&-6\\0&-6&-11\end{pmatrix}
 \longrightarrow
 \begin{pmatrix}1&2&3\\0&-3&-6\\0&0&1\end{pmatrix}.
\]
Thus the determinant is $-3$.
Replacing the last entry $10$ by $9$ gives a zero row after elimination
and determinant zero.
The cyclic permutation matrix
$\begin{pmatrix}0&1&0\\0&0&1\\1&0&0\end{pmatrix}$ has determinant $1$,
because two row swaps turn it into $I_3$.
\end{example}

\subsection{The cofactor inverse formula}

For reference, the larger inverse formula mentioned in
Chapter~\ref{chap:algebra-matrices} has a compact description.
Delete row $i$ and column $j$ to obtain the minor matrix $A_{ij}$, and
define the cofactor $C_{ij}=(-1)^{i+j}\det A_{ij}$.
The \emph{adjugate} is the transpose of the cofactor matrix:
$\operatorname{adj}(A)=(C_{ji})$.
Cofactor expansion gives
\[
 A\operatorname{adj}(A)=\operatorname{adj}(A)A=(\det A)I_n.
\]
Indeed, a diagonal entry is a determinant expansion; an off-diagonal
entry is the expansion of a matrix with two identical rows or columns,
whose determinant is zero.
Consequently,
\[
 \boxed{A^{-1}=\frac{\operatorname{adj}(A)}{\det A}\quad(\det A\ne0).}
\]
For larger matrices, row reduction is usually more practical.
This cofactor explanation supplements the explicit small-matrix formulae
in the source.

\section{The characteristic polynomial}

\begin{definition}
The \emph{characteristic polynomial} of $A$ is
\[
 c_A(t)=\det(tI_n-A).
\]
It is a monic polynomial of degree $n$.
\end{definition}

The convention $\det(tI-A)$ fixes the leading sign.
Some texts use $\det(A-tI)$, which differs by the factor $(-1)^n$
and has the same roots.

\begin{theorem}
$\lambda\in F$ is an eigenvalue of $A$ if and only if $c_A(\lambda)=0$.
\end{theorem}
\begin{proof}
\[
 c_A(\lambda)=0
 \ \Longleftrightarrow\ \lambda I_n-A\text{ is singular}
 \ \Longleftrightarrow\ N(\lambda I_n-A)\ne\{0_n\}.
\]
A nonzero vector in this kernel satisfies $Av=\lambda v$.
\end{proof}

For $A=\begin{pmatrix}a&b\\c&d\end{pmatrix}$,
\[
 c_A(t)=(t-a)(t-d)-bc=t^2-(a+d)t+(ad-bc).
\]
For a triangular matrix,
$c_A(t)=\prod_{j=1}^n(t-a_{jj})$, so its eigenvalues are precisely its
diagonal entries.

\begin{corollary}[Source Corollary 4.21]
Every complex $n\times n$ matrix with $n\geq1$ has an eigenvalue and a
nonzero eigenvector over $\mathbb C$.
\end{corollary}
\begin{proof}
The fundamental theorem of algebra gives a complex root of $c_A$.
The theorem then gives a nonzero vector in the corresponding kernel.
\end{proof}

This conclusion need not hold over $\mathbb R$: the real rotation example
has characteristic polynomial $t^2+1$.
Existence of a complex eigenvector also does not imply existence of a
whole eigenvector basis.

\subsection{Repeated roots}

The \emph{algebraic multiplicity} of an eigenvalue is its multiplicity as
a root of $c_A(t)$.
Its \emph{geometric multiplicity} is $\dim E_\lambda(A)$.
This terminology helps distinguish repeated roots from the number of
independent eigenvectors they supply.
For $J$, $c_J(t)=t^2(t-3)$: the zero eigenvalue has both multiplicities
two, and the eigenvalue $3$ has both multiplicities one.
For $\begin{pmatrix}1&1\\0&1\end{pmatrix}$,
$c_A(t)=(t-1)^2$ but $\dim E_1(A)=1$.

\section{A complete diagonalisation procedure}

\begin{enumerate}
\item Find $c_A(t)$ and its roots in the chosen field $F$.
\item For each distinct root $\lambda$, row-reduce $A-\lambda I_n$ and
      find a basis of $E_\lambda(A)$.
\item Add the eigenspace dimensions. If the total is less than $n$,
      $A$ is not diagonalisable over $F$.
\item Otherwise place the eigenvectors in the columns of $P$ and their
      eigenvalues, in the same order, on the diagonal of $D$.
\item Verify $AP=PD$. This checks the result without explicitly computing
      $P^{-1}$.
\end{enumerate}

\begin{example}[Three distinct eigenvalues]
For
\[
 A=\begin{pmatrix}1&2&3\\0&4&5\\0&0&6\end{pmatrix},
 \qquad c_A(t)=(t-1)(t-4)(t-6).
\]
Solving the three kernel systems gives, with convenient integer scalings,
\[
 E_1=\spanop\left(\begin{pmatrix}1\\0\\0\end{pmatrix}\right),\quad
 E_4=\spanop\left(\begin{pmatrix}2\\3\\0\end{pmatrix}\right),\quad
 E_6=\spanop\left(\begin{pmatrix}16\\25\\10\end{pmatrix}\right).
\]
Thus
\[
 P=\begin{pmatrix}1&2&16\\0&3&25\\0&0&10\end{pmatrix},\qquad
 D=\diag(1,4,6),\qquad AP=PD.
\]
Rescaling any eigenvector by a nonzero scalar changes $P$ but does not
change its eigenvalue.
\end{example}

\begin{example}[A repeated eigenvalue that causes no problem]
For the all-ones matrix,
\[
 P=\begin{pmatrix}-1&-1&1\\1&0&1\\0&1&1\end{pmatrix},\qquad
 D=\diag(0,0,3).
\]
The first two columns form a basis of $E_0(J)$ and the last spans $E_3(J)$.
They form an eigenvector basis, so $P^{-1}JP=D$.
\end{example}

\begin{proposition}[Only one eigenvalue]
Suppose the only eigenvalue of $A$ in $F$ is $\lambda$.
Then $A$ is diagonalisable over $F$ if and only if $A=\lambda I_n$.
\end{proposition}
\begin{proof}
If $A$ is diagonalisable, its diagonal form must be $\lambda I_n$.
Then $A=P(\lambda I_n)P^{-1}=\lambda I_n$.
The converse holds because $\lambda I_n$ is already diagonal.
\end{proof}

\begin{example}
The matrix
\[
 \begin{pmatrix}
 2&2&3&4\\0&2&5&6\\0&0&2&7\\0&0&0&2
 \end{pmatrix}
\]
has characteristic polynomial $(t-2)^4$.
It is not $2I_4$, so it is not diagonalisable.
\end{example}

% !TeX root = ../algebra-1.tex
\chapter{Orthogonality}\label{chap:algebra-orthogonality}

\section{Conjugation, inner products and norms}

Complex conjugation is applied entry by entry to vectors and matrices:
\[
 \overline{A+B}=\overline A+\overline B,\qquad
 \overline{AB}=\overline A\,\overline B,\qquad
 \overline{A^{\mathsf T}}=(\overline A)^{\mathsf T}.
\]
The \emph{conjugate transpose}, or adjoint, is $A^*=\overline A^{\mathsf T}$.
Transposition reverses products, so $(AB)^*=B^*A^*$.
For real matrices, $A^*=A^{\mathsf T}$.

\begin{definition}[The source's inner product convention]
For $u,v\in F^n$, define
\[
 \boxed{u\cdot v=u^{\mathsf T}\overline v
       =\sum_{j=1}^n u_j\overline{v_j}.}
\]
The \emph{norm}, or length, is
\[
 \|u\|=\sqrt{u\cdot u}=\sqrt{|u_1|^2+\cdots+|u_n|^2}.
\]
\end{definition}

For $z=x+iy$, $|z|^2=z\overline z=x^2+y^2$, whereas
$|z|=\sqrt{x^2+y^2}$.
Thus $u\cdot u$ is nonnegative real, and is zero exactly when $u=0_n$.
For real vectors the inner product is the usual dot product.

\begin{proposition}[Inner product rules]
For compatible vectors and matrices,
\begin{align*}
 u\cdot v&=\overline{v\cdot u},\\
 (au+bw)\cdot v&=a(u\cdot v)+b(w\cdot v),\\
 u\cdot(av+bw)&=\overline a(u\cdot v)+\overline b(u\cdot w),\\
 (Au)\cdot v&=u\cdot(A^*v).
\end{align*}
\end{proposition}
\begin{proof}
The first three follow by expanding the coordinate sums. For the last,
\[
 (Au)\cdot v=u^{\mathsf T}A^{\mathsf T}\overline v
 =u^{\mathsf T}\overline{A^*v}=u\cdot(A^*v).
\]
\end{proof}

\begin{remark}
The inner product is linear in the \textbf{first} argument and
conjugate-linear in the second: a scalar leaving the second argument
must be conjugated. Other books may use the opposite convention.
All projection and Gram--Schmidt coefficients below follow our convention.
\end{remark}

\begin{example}[Checking the complex arithmetic]
Let $u=(2+i,-1-i)^{\mathsf T}$, $v=(1,1)^{\mathsf T}$ and
$w=(-1,2)^{\mathsf T}$. Then
\begin{align*}
 u\cdot u&=|2+i|^2+|-1-i|^2=5+2=7,\\
 u\cdot v&=(2+i)+(-1-i)=1,\\
 v\cdot w&=-1+2=1.
\end{align*}
Thus $\|u\|=\sqrt7$. The value of $u\cdot v$ corrects the source's arithmetic.
\end{example}

\section{Orthogonality, Cauchy--Schwarz and angles}

\begin{definition}
Vectors $u,v$ are \emph{orthogonal} if $u\cdot v=0$.
A sequence is orthogonal if each pair of distinct vectors is orthogonal.
\end{definition}

Conjugate symmetry makes the relation symmetric.
For orthogonal $u,v$, expansion gives Pythagoras's identity:
\[
 \|u+v\|^2=(u+v)\cdot(u+v)=\|u\|^2+\|v\|^2.
\]

\begin{theorem}[Cauchy--Schwarz inequality]
For $u,v\in F^n$,
\[
 \boxed{|u\cdot v|\leq\|u\|\,\|v\|.}
\]
Equality holds if and only if $u,v$ are linearly dependent.
\end{theorem}
\begin{proof}
If $v=0_n$, both sides are zero and the pair is dependent.
Otherwise put $c=(u\cdot v)/(v\cdot v)$ and $w=u-cv$.
Then $w\cdot v=0$, so Pythagoras applied to $u=cv+w$ gives
\[
 \|u\|^2=|c|^2\|v\|^2+\|w\|^2
 \geq\frac{|u\cdot v|^2}{\|v\|^2}.
\]
Multiply by $\|v\|^2$ and take square roots.
Equality holds exactly when $w=0$, that is, $u=cv$.
\end{proof}

\begin{definition}[Angles between real vectors]
For nonzero $u,v\in\mathbb R^n$, their angle $\theta\in[0,\pi]$ satisfies
\[
 \cos\theta=\frac{u\cdot v}{\|u\|\,\|v\|}.
\]
\end{definition}

Cauchy--Schwarz ensures the quotient is in $[-1,1]$.
Real nonzero vectors are orthogonal exactly when their angle is $\pi/2$.
For complex vectors the quotient may be nonreal, so this ordinary
real-angle definition cannot simply be used unchanged.
Orthogonality itself remains well-defined over $\mathbb C$.

\begin{proposition}
An orthogonal sequence of nonzero vectors is linearly independent.
\end{proposition}
\begin{proof}
If $\sum_j a_jv_j=0$, take the inner product with $v_k$ on the right.
This gives $a_k\|v_k\|^2=0$. Since $\|v_k\|^2>0$, every $a_k=0$.
\end{proof}

\section{Orthonormal bases}

\begin{definition}
A sequence $q_1,\ldots,q_k$ is \emph{orthonormal} if
\[
 q_i\cdot q_j=\begin{cases}1&i=j,\\0&i\ne j.\end{cases}
\]
Thus its vectors are mutually orthogonal and have norm one.
\end{definition}

An orthogonal nonzero sequence becomes orthonormal by replacing each
$v_j$ with $v_j/\|v_j\|$, which does not change its span.

\begin{theorem}[Coordinates in an orthonormal basis]
If $\mathcal B=(q_1,\ldots,q_d)$ is an orthonormal basis of $U$, then
for $v\in U$,
\[
 \boxed{v=\sum_{j=1}^d(v\cdot q_j)q_j,\qquad
 [v]_{\mathcal B}=(v\cdot q_1,\ldots,v\cdot q_d)^{\mathsf T}.}
\]
\end{theorem}
\begin{proof}
Write $v=\sum_j c_jq_j$ and take the inner product with $q_k$ on the
right. Linearity and orthonormality give $v\cdot q_k=c_k$.
\end{proof}

\begin{example}
An orthonormal basis of $\mathbb R^2$ is
\[
 q_1=\frac1{\sqrt2}(1,1)^{\mathsf T},\qquad
 q_2=\frac1{\sqrt2}(1,-1)^{\mathsf T}.
\]
For $v=(1,3)^{\mathsf T}$, the coordinates are
$v\cdot q_1=2\sqrt2$ and $v\cdot q_2=-\sqrt2$.
Hence $v=2\sqrt2\,q_1-\sqrt2\,q_2$.
\end{example}

\section{Orthogonal projection}

\begin{definition}
The projection of $v$ onto the direction of a nonzero vector $u$ is
\[
 \boxed{\operatorname{proj}_u(v)=\frac{v\cdot u}{u\cdot u}\,u.}
\]
\end{definition}

The projection belongs to $\spanop(u)$, and its residual is orthogonal to $u$:
\[
 \left(v-\frac{v\cdot u}{u\cdot u}u\right)\cdot u=0.
\]
It is the unique vector in $\spanop(u)$ with this property:
if $v-cu$ is orthogonal to $u$, the equation forces
$c=(v\cdot u)/(u\cdot u)$.

\begin{center}
% Drawing only: edit the chapter text in chapters/algebra-05-orthogonality.tex.
\begin{tikzpicture}[>=Stealth,scale=1.05]
 \coordinate (O) at (0,0);
 \coordinate (P) at (3,0);
 \coordinate (V) at (3,1.8);
 \draw[->,muted] (-0.4,0)--(4.4,0) node[right] {$\spanop(u)$};
 \draw[->,thick,ink] (O)--(V) node[midway,above left] {$v$};
 \draw[->,very thick,accent] (O)--(P)
   node[midway,below] {$p=\operatorname{proj}_u(v)$};
 \draw[->,dashed,ink] (P)--(V) node[midway,right] {$v-p$};
 \draw (2.8,0)--(2.8,0.2)--(3,0.2);
 \node[below left] at (O) {$0$};
\end{tikzpicture}

\end{center}

For real vectors with angle $\theta$, the signed component along
$u/\|u\|$ is $\|v\|\cos\theta$.
The projection's length is $\|v\|\,|\cos\theta|$;
the absolute value matters for an obtuse angle.

\subsection{Projection onto a subspace}

For an orthonormal basis $q_1,\ldots,q_d$ of $U$, define
\[
 p=\operatorname{proj}_U(v)=\sum_{j=1}^d(v\cdot q_j)q_j.
\]
Then $p\in U$ and $v-p$ is orthogonal to every vector of $U$.
These properties determine $p$ uniquely, independently of the basis:
if $p,p'$ both work, their difference belongs to $U$ and is orthogonal
to $U$, hence to itself. It has norm zero, so $p=p'$.

\begin{proposition}[Closest point interpretation; additional clarification]
The projection $p$ is the unique vector in $U$ closest to $v$.
\end{proposition}
\begin{proof}
For $y\in U$, $v-y=(v-p)+(p-y)$ is an orthogonal decomposition.
Thus $\|v-y\|^2=\|v-p\|^2+\|p-y\|^2\geq\|v-p\|^2$,
with equality exactly when $y=p$.
\end{proof}

\section{The Gram--Schmidt process}

Given independent $b_1,\ldots,b_k$, subtract projections onto the
previously constructed directions:
\begin{align*}
 v_1&=b_1,\\
 v_2&=b_2-\frac{b_2\cdot v_1}{v_1\cdot v_1}v_1,\\
 v_j&=b_j-\sum_{i=1}^{j-1}
             \frac{b_j\cdot v_i}{v_i\cdot v_i}v_i
             \qquad(2\leq j\leq k).
\end{align*}
Finally set $q_j=v_j/\|v_j\|$.

\begin{theorem}[Gram--Schmidt]
Each $v_j$ is nonzero, the $v_j$ are orthogonal, and
\[
 \spanop(v_1,\ldots,v_j)=\spanop(b_1,\ldots,b_j)
 \quad(1\leq j\leq k).
\]
The $q_j$ are orthonormal. Every subspace of $F^n$ therefore has an
orthonormal basis.
\end{theorem}
\begin{proof}
Induct on $j$. Independence makes $v_1=b_1$ nonzero.
Assume the assertions hold through $j-1$.
For $\ell<j$, take the inner product with $v_\ell$ on the right:
\[
 v_j\cdot v_\ell=b_j\cdot v_\ell
 -\frac{b_j\cdot v_\ell}{v_\ell\cdot v_\ell}
 (v_\ell\cdot v_\ell)=0.
\]
All other terms vanish by orthogonality.
The definition subtracts a vector in the preceding span, so adjoining
$v_j$ or adjoining $b_j$ gives the same enlarged span.
If $v_j=0$, then $b_j$ belongs to the span of its predecessors, contrary
to independence. All denominators are therefore positive and
normalisation is valid. Starting with a basis of a subspace proves the
last assertion.
\end{proof}

\begin{remark}
Over $\mathbb C$, the numerator is $b_j\cdot v_i$, not $v_i\cdot b_j$,
under our convention. They are conjugates and need not be equal.
\end{remark}

\begin{example}[A complete three-dimensional calculation]
Start with
\[
 b_1=\begin{pmatrix}1\\2\\3\end{pmatrix},\quad
 b_2=\begin{pmatrix}4\\5\\6\end{pmatrix},\quad
 b_3=\begin{pmatrix}7\\8\\10\end{pmatrix}.
\]
Their determinant is $-3$, so they are independent.
As $\|b_1\|^2=14$ and $b_2\cdot b_1=32$,
\[
 v_1=b_1,\qquad
 v_2=b_2-\frac{16}{7}b_1
     =\frac37\begin{pmatrix}4\\1\\-2\end{pmatrix}.
\]
Projection onto a direction is unchanged by rescaling its generator,
so use $w_2=(4,1,-2)^{\mathsf T}$ at the next step.
We have $b_3\cdot b_1=53$, $b_3\cdot w_2=16$ and $\|w_2\|^2=21$.
Thus
\begin{align*}
 v_3&=b_3-\frac{53}{14}b_1-\frac{16}{21}w_2\\
 &=\frac16\begin{pmatrix}1\\-2\\1\end{pmatrix}.
\end{align*}
After normalisation,
\[
 \boxed{
 q_1=\frac1{\sqrt{14}}\begin{pmatrix}1\\2\\3\end{pmatrix},\quad
 q_2=\frac1{\sqrt{21}}\begin{pmatrix}4\\1\\-2\end{pmatrix},\quad
 q_3=\frac1{\sqrt6}\begin{pmatrix}1\\-2\\1\end{pmatrix}.}
\]
Their pairwise dot products vanish and each has norm one.
\end{example}

\method{Keep the orthogonal vectors unnormalised during subtraction,
then normalise at the end to avoid unnecessary square roots.
A zero residual indicates dependence in the input sequence.}

\section{Orthogonal and unitary matrices}

\begin{definition}
A real square matrix $Q$ is \emph{orthogonal} if $Q^{\mathsf T}Q=I_n$.
A complex square matrix $Q$ is \emph{unitary} if $Q^*Q=I_n$.
\end{definition}

For an orthogonal matrix, $Q^{-1}=Q^{\mathsf T}$; for a unitary matrix,
$Q^{-1}=Q^*$.
A square matrix with a left inverse is invertible, by the
invertibility theorem.

\begin{proposition}
A real square matrix is orthogonal exactly when its columns are an
orthonormal basis of $\mathbb R^n$.
The complex analogue holds with unitary matrices and $\mathbb C^n$.
\end{proposition}
\begin{proof}
For real columns $q_i$, entry $(i,j)$ of $Q^{\mathsf T}Q$ is $q_i\cdot q_j$.
For complex columns, entry $(i,j)$ of $Q^*Q$ is $q_j\cdot q_i$ under
our convention. It equals $\delta_{ij}$ exactly when the columns are
orthonormal. There are $n$ independent columns, so they form a basis.
\end{proof}

Orthogonal and unitary matrices preserve inner products:
\[
 (Qu)\cdot(Qv)=u\cdot(Q^*Qv)=u\cdot v.
\]
Thus they preserve lengths and orthogonality.
Real orthogonal matrices also preserve real angles.
Taking determinants in $Q^{\mathsf T}Q=I_n$ gives
$(\det Q)^2=1$, so $\det Q=1$ or $-1$.

\begin{remark}
An orthogonal \emph{matrix} has orthonormal columns, including length one.
Merely having orthogonal nonzero columns is insufficient:
$2I_n$ has orthogonal columns but is not an orthogonal matrix.
\end{remark}

\section{Real symmetric matrices and the spectral theorem}

\begin{definition}
A real matrix is \emph{symmetric} if $A^{\mathsf T}=A$.
A complex matrix is \emph{Hermitian} if $A^*=A$.
\end{definition}

\begin{lemma}
Every complex eigenvalue of a real symmetric matrix is real.
Eigenvectors for distinct eigenvalues are orthogonal.
\end{lemma}
\begin{proof}
Let $Av=\lambda v$ with $v\in\mathbb C^n$ nonzero. Since $A^*=A$,
\[
 \lambda\|v\|^2=(Av)\cdot v
 =v\cdot(Av)=\overline\lambda\|v\|^2.
\]
The positive norm squared gives $\lambda=\overline\lambda$.
If $Au=\lambda u$ and $Av=\mu v$, then
\[
 \lambda(u\cdot v)=(Au)\cdot v=u\cdot(Av)=\mu(u\cdot v).
\]
Distinct real $\lambda,\mu$ force $u\cdot v=0$.
\end{proof}

\begin{lemma}[A real eigenvector can be chosen]
A real symmetric matrix of positive size has a nonzero real eigenvector.
\end{lemma}
\begin{proof}
Its characteristic polynomial has a complex root $\lambda$; the
previous lemma makes $\lambda$ real.
Write a complex eigenvector as $v=x+iy$ with $x,y\in\mathbb R^n$.
Because $A,\lambda$ are real, $Av=\lambda v$ gives
$Ax=\lambda x$ and $Ay=\lambda y$.
At least one of $x,y$ is nonzero, giving a real eigenvector.
\end{proof}

\begin{theorem}[Real spectral theorem; source Theorem 5.22, strengthened form]
Every real symmetric matrix $A$ has an orthonormal basis of real
eigenvectors. Equivalently, a real orthogonal matrix $Q$ exists with
\[
 \boxed{Q^{\mathsf T}AQ=D,\qquad A=QDQ^{\mathsf T},}
\]
where $D$ is real and diagonal. In particular, $A$ is diagonalisable
over $\mathbb R$.
\end{theorem}
\begin{proof}
Induct on $n$, with $n=1$ immediate.
Choose a real unit eigenvector $q_1$ with eigenvalue $\lambda$.
Extend it to a basis and apply Gram--Schmidt.
Since $q_1$ is already a unit vector, it remains the first vector.
This gives an orthogonal $Q_0=[\,q_1\ q_2\ \cdots\ q_n\,]$.

The first column of $Q_0^{\mathsf T}AQ_0$ is
$(\lambda,0,\ldots,0)^{\mathsf T}$.
Symmetry forces the same zeros in the first row, so
\[
 Q_0^{\mathsf T}AQ_0=\begin{pmatrix}\lambda&0\\0&B\end{pmatrix},
\]
where $B$ is real symmetric of size $n-1$.
By induction, an orthogonal $Q_1$ makes $Q_1^{\mathsf T}BQ_1=D_1$ diagonal.
Then
\[
 Q=Q_0\begin{pmatrix}1&0\\0&Q_1\end{pmatrix}
\]
is orthogonal and $Q^{\mathsf T}AQ=\diag(\lambda,D_1)$.
Its columns are the required eigenvector basis.
\end{proof}

\begin{example}
For $A=\begin{pmatrix}2&1\\1&2\end{pmatrix}$, take
$Q=\frac1{\sqrt2}\begin{pmatrix}1&1\\1&-1\end{pmatrix}$.
Its columns are orthonormal eigenvectors and
$Q^{\mathsf T}AQ=\diag(3,1)$.
This is the diagonalisation from Chapter~\ref{chap:algebra-linear-maps},
now with an orthogonal change of basis.
\end{example}

\begin{remark}[Complex symmetric is different from Hermitian]
Over $\mathbb C$, $A^{\mathsf T}=A$ alone does not imply diagonalisability.
For example,
\[
 A=\begin{pmatrix}1&i\\i&-1\end{pmatrix}
\]
is symmetric and nonzero but $A^2=0$.
Any eigenvalue must be zero; diagonalisation would therefore force
$A=0$, a contradiction.
The complex spectral theorem requires $A^*=A$:
a Hermitian matrix has real eigenvalues and a unitary diagonalisation
$Q^*AQ=D$. This is the Hermitian analogue mentioned in the source.
\end{remark}

\end{document}
